(1) 由第一定律 \(dU=T\,dS-p\,dV\)，对 \(U=aVT^4\) 微分得到 \(dU=4aVT^3dT+aT^4dV\)。又 \(p=aT^4/3\)，因此 \(T\,dS=4aVT^3dT+(4aT^4/3)dV\)，积分得 \(S=(4a/3)VT^3+S_0\)；取零温熵零定常数。
(2) 可逆绝热要求 \(S\) 常量，所以 \(VT^3=\mathrm{const}\)，等价于 \(TV^{1/3}=\mathrm{const}\)。该关系反映辐射体的熵随体积与温度的三次幂增长，不等同于固定粒子数分子气体的泊松绝热方程。
(3) 焓定义为 \(H=U+pV\)，故 \(H=aVT^4+aVT^4/3=4aVT^4/3\)。定容热容 \(C_V=(\partial U/\partial T)_V=4aVT^3\)；它随温度和容积改变，并非常量，因此不能直接应用恒热容理想气体公式。
(4) 由 \(T\propto V^{-1/3}\) 与 \(p=aT^4/3\) 联立，得到 \(p\propto V^{-4/3}\)，故 \(pV^{4/3}=\mathrm{const}\)，有效绝热指数为 \(\gamma_{\rm rad}=4/3\)。这是平衡辐射气体的状态规律，而非单原子理想气体的 \(5/3\)。
